\section{Background and Related Work}\label{sec:related}

\subsection{The Fuchs et al. 2015 concept}

Fuchs, Langer and Trinitis~\cite{fuchs2015ftrfs} proposed the
original FTRFS concept at ARCS~2015 as a research contribution
exploring fault-tolerance for filesystems exposed to radiation
environments. The 2015 work established the principle of integrating
forward error correction at the filesystem layer rather than
delegating it to lower storage tiers. It did not target the Linux
mainline kernel as we know it today: the kernel infrastructure
available in 2015 (\texttt{lib/reed\_solomon} primarily used by
MTD/NAND), the validation tooling (xfstests scope, syzkaller maturity,
KASAN), and the surrounding standards (MIL-STD-882E adoption for
software, NIST SP~800-193~\cite{nistsp800193}, NIS2 EU directive,
French OIV regulation) have all evolved substantially. The present
work is a 2026 implementation responding to the 2015 problem
statement under contemporary constraints.

\subsection{Linux/BSD filesystem landscape (2026)}

\Cref{tab:fs-comparison} summarizes how FTRFS positions against
the production filesystems available to a Linux deployment as of
April~2026. The comparison restricts itself to verifiable facts
from each project's documentation and source tree; it does not
attempt to rank filesystems by overall merit or suitability for
non-radiation workloads.

\begin{table*}[t]
\centering
\small
\caption{Filesystem metadata protection landscape (April 2026).}
\label{tab:fs-comparison}
\begin{tabularx}{\textwidth}{lXXXXl}
\toprule
\textbf{FS} &
\textbf{Metadata checksum} &
\textbf{Metadata FEC} &
\textbf{On-disk event journal} &
\textbf{Fail-closed} &
\textbf{Tree} \\
\midrule
ext4
  & CRC32C (optional)
  & none
  & none
  & best-effort
  & mainline \\
btrfs
  & CRC32C / xxhash / sha256
  & none
  & \texttt{dev\_stats} (counters)
  & best-effort
  & mainline \\
bcachefs
  & CRC32C / xxhash (optional)
  & none
  & none
  & best-effort
  & mainline \\
ZFS (OpenZFS)
  & fletcher4 / sha256
  & RAID-Z (data only)
  & in-memory pool log
  & partial
  & out-of-tree \\
\textbf{FTRFS}
  & \textbf{CRC32 + RS FEC layered}
  & \textbf{RS(255,239) on inode, bitmap, superblock}
  & \textbf{64-entry ring buffer in SB, CRC-protected}
  & \textbf{strict (version, scheme, features)}
  & out-of-tree (RFC v3) \\
\bottomrule
\end{tabularx}
\end{table*}

To our knowledge, no other Linux filesystem combines (i) universal
Reed--Solomon FEC on metadata structures, (ii) a persistent on-disk
journal of correctable events suitable for forensic post-mortem,
and (iii) strict fail-closed semantics on multiple structural axes
simultaneously. Individual properties exist in isolation: ZFS has
RAID-Z FEC on data blocks; ext4 has metadata CRC; btrfs records
device-level error counters. Their conjunction in a single
metadata trust chain is, as far as we have surveyed, specific to
FTRFS.

\subsection{Community engagement on linux-fsdevel}

The project has been submitted to the linux-fsdevel mailing list
in three RFC iterations (v1, v2, v3) between January and April
2026. The threads are publicly archived on
\url{lore.kernel.org} and have engaged senior reviewers active in
filesystem maintenance. The engagement has produced substantial
design feedback that has shaped the architecture documented here ---
notably the layered CRC32+RS defense, the universal protection
scheme replacing the original per-inode opt-in, and the strict
fail-closed posture on version mismatches. We treat the present
report as a structured complement to those mailing-list threads,
not a replacement.

